\documentclass{article}
\usepackage{amsmath}
\usepackage{graphicx}
\usepackage{hyperref}

\title{Checkpoint 1 Report: Negative-Weight Single-Source Shortest Paths in Near-Linear Time}
\author{Team: 68 \\
Muhammad Daniyal Farooqui, Jibran Sheikh}

\begin{document}

\maketitle

\section{Paper Selection Proposal}

\subsection{Paper Details}

\begin{itemize}
    \item \textbf{Title:} Negative-Weight Single-Source Shortest Paths in Near-Linear Time
    \item \textbf{Authors:} Aaron Bernstein, Danupon Nanongkai, Christian Wulff-Nilsen
    \item \textbf{Conference:} Communications of the ACM Vol. 68, No. 2
    \item \textbf{Year:} 2023 (updated)
    \item \textbf{Link:} \url{https://dl.acm.org/doi/10.1145/3631536}
\end{itemize}

\subsection{Summary}

The paper addresses the Single-Source Shortest Paths (SSSP) problem in directed graphs with negative edge weights, a fundamental challenge in graph algorithms. Traditional methods like Bellman-Ford have high time complexity \( O(mn) \), which is impractical for large graphs. Recent approaches rely on complex optimization techniques, introducing significant overhead. 

The authors present a **combinatorial algorithm** achieving near-linear runtime \( O(m \log^6 n \log W) \), significantly simplifying prior methods. 

Key innovations include:
\begin{itemize}
    \item **Price functions** to transform edge weights into non-negative values.
    \item **Low-diameter decomposition** to partition the graph into manageable components.
    \item **Recursive scaling** to progressively eliminate negative weights.
\end{itemize}

\subsection{Justification}

The paper was selected for the following reasons:

\begin{itemize}
    \item \textbf{Theoretical Impact:} Resolves a decades-old problem with a near-optimal solution.
    \item \textbf{Practical Relevance:} Applicable to logistics (e.g., routes with fuel rebates), financial arbitrage, and network optimization.
    \item \textbf{Simplicity:} Avoids complex continuous optimization, making it more accessible for implementation.
    \item \textbf{Classroom Relevance:} We have already studied both algorithms in class, and it would be interesting to work on the shortcomings of those algorithms and overcome them using a single unified algorithm.
\end{itemize}

\subsection{Implementation Feasibility}

\textbf{Resources:}
\begin{itemize}
    \item \textbf{Theoretical Foundation:} Detailed pseudocode and proofs are provided.
    \item \textbf{Code Repositories:} No public code exists, but standard graph libraries (e.g., NetworkX, Boost Graph) can be adapted.
    \item \textbf{Datasets:} Synthetic graphs with negative weights.
\end{itemize}

\textbf{Theoretical Descriptions:}
\begin{itemize}
    \item Efficient handling of probabilistic graph decomposition.
    \item Optimizing priority queues for the hybrid Dijkstra-Bellman-Ford.
\end{itemize}

\subsection{Team Responsibilities}
\begin{itemize}
    \item \textbf{Reading and Analyzing the Paper:} Muhammad Daniyal Farooqui
    \item \textbf{Coding and Implementation:} Jibran Sheikh, Muhammad Daniyal Farooqui
    \item \textbf{Writing:} Jibran Sheikh
\end{itemize}

\end{document}
